% Figure 2: why answer-span labels and relevance are proxies (illustrative values, not measurements).
% Emphasis form: the utility column is the point (blue), the two proxies are context (gray / marks).
\begin{tikzpicture}[
    font=\scriptsize,
    para/.style={anchor=north west, text width=3.05cm, align=left, inner sep=0pt, text=inkP},
    head/.style={font=\scriptsize\bfseries, text=inkS, align=center, anchor=south},
  ]
  \def\rowh{9.5mm}
  \def\barh{2.4mm}
  \def\barmax{1.05cm}
  % column x positions
  \def\xspan{3.75cm}
  \def\xrel{4.35cm}
  \def\xbeta{5.95cm}

  \node[anchor=south west, text width=7.3cm, align=left, inner sep=0pt, font=\footnotesize, text=inkP]
    at (0, 0.8cm) {\textbf{Câu hỏi:} Đạo diễn của phim \emph{Mùi đu đủ xanh} sinh năm nào?\quad
    \textbf{Đáp án:} 1962};
  \node[head] at (\xspan, 0.05cm) {Chứa\\ đáp án};
  \node[head] at ($(\xrel,0.05cm)+(0.5*\barmax,0)$) {Độ liên quan\\ (reranker)};
  \node[head] at ($(\xbeta,0.05cm)+(0.5*\barmax,0)$) {Mức đóng góp\\ $\beta$ (của chúng tôi)};
  \draw[axisc, line width=0.4pt] (0,0) -- (7.3cm,0);

  % rows: text / contains answer / relevance / beta
  \foreach \i/\txt/\hasans/\rel/\bval in {
      1/{\textbf{P1} (cầu nối). \emph{Mùi đu đủ xanh} (1993) là phim của đạo diễn Trần Anh Hùng.}/0/0.85/0.90,
      2/{\textbf{P2}. Trần Anh Hùng sinh năm 1962, là đạo diễn người Pháp gốc Việt.}/1/0.40/0.95,
      3/{\textbf{P3}. \emph{Mùi đu đủ xanh} đoạt giải Caméra d'Or tại Cannes năm 1993.}/0/0.90/0.05,
      4/{\textbf{P4}. Đà Lạt là thành phố thuộc tỉnh Lâm Đồng.}/0/0.05/0.02} {
    \pgfmathsetlengthmacro{\ytop}{-1.5mm-(\i-1)*\rowh}
    \pgfmathsetlengthmacro{\ymid}{\ytop-3.2mm}
    \node[para] at (0, \ytop) {\txt};
    \ifnum\hasans=1
      \node[text=inkP, font=\small] at (\xspan, \ymid) {\checkmark};
    \else
      \node[text=inkM] at (\xspan, \ymid) {--};
    \fi
    % baselines, then bars (square at the baseline, rounded data end)
    \draw[axisc, line width=0.4pt] (\xrel, \ymid-2mm) -- (\xrel, \ymid+2mm);
    \draw[axisc, line width=0.4pt] (\xbeta, \ymid-2mm) -- (\xbeta, \ymid+2mm);
    \fill[inkM, rounded corners=0.7pt] (\xrel, \ymid-0.5*\barh) rectangle ++(\rel*\barmax, \barh);
    \fill[vzblue, rounded corners=0.7pt] (\xbeta, \ymid-0.5*\barh) rectangle ++(\bval*\barmax, \barh);
  }
\end{tikzpicture}
